\documentclass[11pt]{article}
\usepackage{mathblatt}
\usepackage{multicol}


\begin{document}
\pfheftstil
\fancyfoot[C]{\parbox[t]{\textwidth}{\mbfhbox\par\vspace{1.5mm}{\footnotesize\color{mbgrau}{\scriptsize UTU}\hfill\thepage}}}
\pfheftkopf{Satz des Pythagoras \textperiodcentered{} Lösungen}{}
{\small\noindent\textbf{Typische Fehler}\par\smallskip\noindent \textbullet\ Quadrate addiert, obwohl eine Kathete gesucht ist.\par\noindent \textbullet\ falsche Seite als Hypotenuse (längere Kathete; x² = y² + z²).\par\noindent \textbullet\ Wurzel vergessen, c² als Ergebnis.\par\noindent \textbullet\ c = a + b; Wurzel aus jedem Summanden einzeln.\par\noindent \textbullet\ in Figuren: ganze statt halbe Seite, Durchmesser statt Radius, Schräge als Höhe.\par\noindent \textbullet\ Buchstabenfixierung „c ist immer die Hypotenuse“.\par}\medskip
\begin{pfloesung}{}
\lz{1.}{a) H = $r$; b) H = $g$; c) kein rechter Winkel; d) H = $v$}{}{}
\end{pfloesung}
\begin{pfloesung}{}
\lz{2.}{a) $b^2 = a^2 + c^2$; b) $x^2 = y^2 + z^2$; c) $q^2 = p^2 + s^2$}{}{}
\end{pfloesung}
\begin{pfloesung}{}
\lz{3.}{z² = x² + y²}{}{}
\end{pfloesung}
\begin{pfloesung}{}
\lz{4.}{z = √(x² + y²)}{}{}
\end{pfloesung}
\begin{pfloesung}{}
\lz{5.}{w = √(u² + v²)}{}{}
\end{pfloesung}
\begin{pfloesung}{}
\lz{6.}{z = √(x² $-$ y²)}{z² = x² $-$ y²}{}
\end{pfloesung}
\begin{pfloesung}{}
\lz{7.}{Die Summe der Flächeninhalte der Kathetenquadrate ist gleich dem Flächeninhalt des Hypotenusenquadrates}{a² + b² = c²}{}
\end{pfloesung}
\begin{pfloesung}{}
\lz{8.}{$c = 39$ cm}{Gleichung: $c^2 = 36^2 + 15^2$; Zwischenergebnis: $c^2 = 1\,296 + 225 = 1\,521$; Wurzel: $c = \sqrt{1\,521} = 39$ cm}{}
\end{pfloesung}
\begin{pfloesung}{}
\lz{9.}{$\approx 9{,}8$ cm}{$c^2 = 97$}{}
\end{pfloesung}
\begin{pfloesung}{}
\lz{10.}{$\overline{AB}$ = 10√13 $\approx$ 36,1 m}{$\overline{AB}^2$ = 12² + 34² = 1300}{}
\end{pfloesung}
\begin{pfloesung}{}
\lz{11.}{x $\approx$ 170,8 cm}{x² = 170² + 16² = 29 156}{}
\end{pfloesung}
\begin{pfloesung}{}
\lz{12.}{minus}{Kathete gesucht}{}
\end{pfloesung}
\begin{pfloesung}{}
\lz{13.}{$a = 84$ cm}{Gleichung umstellen: $a^2 = 85^2 - 13^2$; Zwischenergebnis: $a^2 = 7\,225 - 169 = 7\,056$; Wurzel: $a = \sqrt{7\,056} = 84$ cm}{}
\end{pfloesung}
\begin{pfloesung}{}
\lz{14.}{$\approx 9{,}2$ cm}{$a^2 = 85$}{}
\end{pfloesung}
\begin{pfloesung}{}
\lz{15.}{$\overline{BC}$ = √65 $\approx$ 8,06 m}{$\overline{BC}$² = 9² $-$ 4² = 65}{}
\end{pfloesung}
\begin{pfloesung}{}
\lz{16.}{$\overline{FA}$ $\approx$ 287,1 m}{$\overline{FA}$² = 384² $-$ 255² = 82 431}{}
\end{pfloesung}
\begin{pfloesung}{}
\lz{17.}{h = √855 $\approx$ 29,2 cm}{h² = 32² $-$ 13² = 855}{}
\end{pfloesung}
\begin{pfloesung}{}
\lz{18.}{$\overline{AD}$ = √144,05 $\approx$ 12,0 m}{$\overline{AD}$² = 14,1² $-$ 7,4² = 144,05}{}
\end{pfloesung}
\begin{pfloesung}{}
\lz{19.}{$\overline{AD}$ = √553 719 $\approx$ 744 m}{$\overline{AD}^2$ = 920² $-$ 541²}{}
\end{pfloesung}
\begin{pfloesung}{}
\lz{20.}{Katheten: die Höhe und der Überstand; Hypotenuse: der Schenkel}{}{}
\end{pfloesung}
\begin{pfloesung}{}
\lz{21.}{$h = 63$ cm}{halbe Grundseite: $\tfrac{32}{2} = 16$ cm; Gleichung umstellen: $h^2 = 65^2 - 16^2$; Zwischenergebnis: $h^2 = 4\,225 - 256 = 3\,969$; Wurzel: $h = \sqrt{3\,969} = 63$ cm}{}
\end{pfloesung}
\begin{pfloesung}{}
\lz{22.}{Gesamthöhe $\approx$ 32,2 m}{Kegelhöhe² = 8,6² $-$ 4,7² = 51,87; Kegelhöhe $\approx$ 7,20 m}{}
\end{pfloesung}
\begin{pfloesung}{}
\lz{23.}{$\overline{BC}$ = √20 $\approx$ 4,47}{$\overline{AB}$ = 2, $\overline{AC}$ = 4; $\overline{BC}$² = 2² + 4² = 20}{}
\end{pfloesung}
\begin{pfloesung}{}
\lz{24.}{$h_s$ $\approx$ 0,79 m}{halbe Grundkante 0,40 m; $h_s$² = 0,40² + 0,68² = 0,6224}{}
\end{pfloesung}
\begin{pfloesung}{}
\lz{25.}{$\overline{AB}$ $\approx$ 60,4 cm}{$\overline{AF}$ = 25 cm, $\overline{BF}$ = 80 $-$ 25 = 55 cm; $\overline{AB}$² = 25² + 55² = 3650}{}
\end{pfloesung}
\begin{pfloesung}{}
\lz{26.}{Radius 3,2 m und Dachhöhe h als Katheten, Dachbalken als Hypotenuse: s = √(3,2² + h²)}{}{}
\end{pfloesung}
\begin{pfloesung}{}
\lz{27.}{√89,96 + 2 $\approx$ 11,5 cm}{Diagonale √(7² + 6,4²) = √89,96 $\approx$ 9,5 cm; Durchmesser 6,4 cm}{}
\end{pfloesung}
\end{document}